\section{What the exposure costs}
\label{sec:measurement}

Section~\ref{sec:obstruction} says the guarantee is lost. It does not say how
often the loss is felt, and that question decides whether the result is a caution
or a curiosity. We measure it on a shipped solver, \texttt{cvx-quadprog}
version $0.4.2$~\cite{cvxquadprog}, whose fast path is the exchange of
Section~\ref{sec:exchange} guarded by the certificate of
Section~\ref{sec:certificate}, with the classical dual walk as its fallback. The
experiment instruments the shipped routines by wrapping them rather than
reimplementing them, so what is measured is what runs. It covers
$\qpPdasFamilies{}$ constraint families, $\qpPdasInstances{}$ instances each, at
$n \in \{\qpPdasSizes\}$, with $G = BB\T + nI$ for $B$ Gaussian and every problem
drawn from a seed derived arithmetically from family, size and instance index.

\begin{figure}[ht]
\centering
\includegraphics[width=0.86\textwidth]{graphs/quadprog_pdas.pdf}
\caption{Left: the fraction of guesses the certificate accepts. Right: set repairs
to convergence. Four families are certified without exception at every size and
settle in \qpPdasRepairsLo{} to \qpPdasRepairsHi{} repairs while the exact walk's
outer iterations grow linearly in $n$. The fifth is constructed so that a guessed
set can be linearly dependent, and it fails.}
\label{fig:pdas}
\end{figure}

\paragraph{The method works, and works well, where the rank condition holds.} On
the four well-posed families (bounds, budget-plus-bounds, dense $C$, and $C$ with
equalities) every guess is certified at every size, and the repair count rises
only from \qpPdasRepairsLo{} to \qpPdasRepairsHi{} across a sixteen-fold range of
$n$, never exceeding \qpPdasRepairsMax{}. The comparison that matters is against
the exact walk on the same instances: on budget-plus-bounds it takes
\qpPdasBudgetOuter{} outer iterations at $n = \qpPdasRefN{}$, and its count grows
linearly where the repair count is very nearly flat. This is why the exchange is
worth having, and it is the reason the negative result matters rather than a
qualification of it.

\paragraph{The rank-deficient guess is real, and invisible by accident.} Over the
$600$ instances of those four families, not one guess produced a rank-deficient
working-set system, from which a reader might conclude the exposure is
theoretical. It is not; it is invisible in random data, where a dense $C$ almost
never has a dependent subset that the guess happens to name. The fifth family
makes it reachable in the most elementary way, by repeating columns of $C$, which
is the construction of Proposition~\ref{prop:reachable}. There the rank test fires
on \qpPdasSingularDup\% of attempts and the certified fraction collapses from
\qpPdasCertSmall\% at $n = \qpPdasNsmall{}$ to \qpPdasCertBig\% at
$n = \qpPdasNbig{}$, against \qpPdasSingularBox\% on the others. The two numbers
bracket the practical question: a benchmark suite of random instances will report
that the fast path never fails, and a model assembled from overlapping rules will
find that it essentially never succeeds.

\paragraph{Where the failures are caught.} Of the candidates reaching the
certificate, \qpPdasAccepted{} were accepted with a worst KKT
residual~\eqref{eq:certificate} of $\qpPdasAcceptWorst$, and \qpPdasRejected{}
were rejected. We report that asymmetry instead of smoothing it: on these
families the certificate never had to reject a converged candidate, because the
failures were caught earlier and more cheaply by the Cholesky factorisation
refusing a dependent guess. The certificate remains the guarantee, and it is what
one must keep, since nothing here bounds how a rank-deficient guess would behave
under a solver that tolerated it. On this evidence, though, the rank test does the
work in practice.

\paragraph{What this does not establish.} The experiments are dense and of modest
size, with $G = BB\T + nI$, so they say nothing about conditioning; the residuals
in Remark~\ref{rem:tolerances} would grow with $\kappa(G)$. That choice also works
against the conclusion drawn here rather than for it, since a well-conditioned $G$
spreads the minimiser so that few bounds bind, where a realistic covariance
concentrates it and binds many. Nor do they establish that duplication is the
common case in practice; they establish that when it occurs the failure rate is
not small.
